\chapter[Conclusão]{Conclusão}

Este é o modelo canônico do curso de Ciência da Computação da UFSJ, use-o sempre para fazer sua monografia.
Ele é baseado no abnTeX2, que é um conjunto de macros LaTeX para formatação de documentos seguindo as normas da ABNT.
O abnTeX2 é mantido por Lauro César Araujo e outros colaboradores.

Agradecimentos também ao projeto UCSMONOGRAPH, disponível em \url{https://github.com/henriquebaron/ucsmonograph}, que trouxe soluções elegantes para simplificar o arquivo .tex base (acredite!) trazendo várias definições diretamente para a classe.

Para utilizar este modelo em editores online, como o Overleaf, você deve colocar o arquivo \texttt{ufsj-abntex2.cls} na mesma pasta.
Use o arquivo \texttt{main.tex} como base, irá adiantar seu trabalho.

No Linux, instale os pacotes \texttt{texlive}, \texttt{texlive-latex-extra}, \texttt{texlive-science}, \texttt{texlive-lang-portuguese} e \texttt{biber} para compilar corretamente. O pacote já contem um Makefile para adiantar. Edite as variáveis conforme necessário.